\documentclass[11pt]{article}
\usepackage[margin=1in]{geometry}
\usepackage{amsmath,amssymb,amsthm,booktabs,hyperref,microtype,array}
\hypersetup{colorlinks=true,linkcolor=blue,pdftitle={Basin Entry and OS Frontier after H1},pdfauthor={Yang--Mills representational dynamics research programme}}
\newtheorem{proposition}{Proposition}
\newtheorem{warning}{Warning}
\title{Basin Entry and the Osterwalder--Schrader Frontier\\Target 58 of the H1 Downstream Programme}
\author{Yang--Mills representational dynamics research programme}
\date{21 July 2026}
\begin{document}
\maketitle
\begin{abstract}
We audit the complete five-target run against physical basin entry and the
Clay mass-gap target.  A positive temporal-transfer countermodel proves that a
fixed one-scale spatial defect ceiling is compatible with arbitrarily small
temporal gaps.  A cutoff family further shows that positive regulated gaps may
vanish in physical units.  The finite programme reaches C2, with C3
conditional bridges; the first physical nonlinear kernel inequality remains
open.
\end{abstract}

\section{What the five targets prove}
The run establishes: (i) exact finite normalized disintegration; (ii) a finite
Harnack criterion for inverse-vacuum moments; (iii) robust multichannel Perron
envelopes; (iv) changing-state-space weight cocycles and the correct
physical-length criterion; and (v) the present frontier audit.  None of the
finite witnesses is asserted to be the complete nonlinear $SU(2)$ measure.

\section{Temporal independence countermodel}
For $\varepsilon>0$, define
\[
 T_\varepsilon=\frac12
 \begin{pmatrix}
 1+e^{-\varepsilon}&1-e^{-\varepsilon}\\
 1-e^{-\varepsilon}&1+e^{-\varepsilon}
 \end{pmatrix}.
\]
It is symmetric, entrywise positive, and has eigenvalues $1$ and
$e^{-\varepsilon}$.
\begin{proposition}
The transfer gap at unit time is $\varepsilon$.  Therefore the same spatial
weighted defect ceiling
\[
 q=0.002430277501841998
\]
is compatible with temporal gaps tending to zero.
\end{proposition}
\begin{proof}
The transfer-matrix gap is
$-\log(e^{-\varepsilon}/1)=\varepsilon$.  The spatial defect kernel is an
independent tensor factor and is unchanged as $\varepsilon\downarrow0$.
\end{proof}

There is also a cutoff version.  Let the temporal lattice spacing be $a$ and
choose $\varepsilon(a)=a^2$.  Every regulator has a positive transfer gap, but
the gap in physical units is
\[
 \frac{\varepsilon(a)}a=a\longrightarrow0.
\]
At $a=2^{-12}$ the regulated physical gap is $2^{-12}=0.000244140625$.

\section{Gate audit}
\begin{center}
\begin{tabular}{p{0.62\linewidth}c}
\toprule
Obligation & Status\\
\midrule
Complete normalized nonlinear $SU(2)$ retained kernel & open\\
Uniform vacuum/Lyapunov weight and root moment & open\\
Complete multichannel weighted ceiling below threshold & open\\
Repeated contraction in accumulated physical length & open\\
Continuum OS construction and reflection positivity & open\\
Dense state-separating temporal observable core & open\\
Identity with nontrivial pure Yang--Mills & open\\
Every compact simple gauge group & open\\
\bottomrule
\end{tabular}
\end{center}
The first failure is the physical normalized kernel.  The current claim level
is C2, not C4.  The conditional statement that dense-core temporal decay in an
already constructed OS theory excludes spectrum in $(0,m)$ is a C3 bridge,
not a construction theorem.

\section{Limit order and next theorem}
The defensible order is: construct finite normalized rows; prove uniform
thermodynamic control; prove cutoff-uniform physical-length contraction;
construct and identify the OS continuum theory; prove density of the decaying
observable core; then apply the spectral-measure implication.

The first open physical target is
\[
 \sup_{a,L,h}\frac{(Q_{a,L}V_{a,L})(h)}{V_{a,L}(h)}<q_{\rm crit},
 \qquad
 \sup_{a,L}\nu_{a,L}(V_{a,L})<\infty,
\]
for every positive-mass history and a weight defined without clustering,
confinement, or a pre-existing gap.

\begin{warning}
A one-scale spatial certificate, finite-volume transfer gap, formal RG flow,
or assumed confinement scale does not establish the Hamiltonian mass
gap.  Continuum existence, OS axioms, observable density, theory identity, and
physical-unit uniformity are separate obligations.
\end{warning}

\section{Conclusion}
The five targets improve the representation calculus and eliminate two major
shortcuts.  They do not solve Yang--Mills.  Eleven audit tests verify the
temporal spectra, cutoff collapse, claim-level gates, and conditional OS
handoff.
\end{document}
